\chapter{Réglages d'émulation}\label{chap:emu-settings}

\section{SID}
\capture{settings-sid}{Réglages SID dans le panneau en direct}
\begin{description}
  \item[Moteur] reSID ou reSIDfp.
  \item[Modèle] 6581 ou 8580 pour les morceaux qui ne le précisent pas, ou pour tous si le choix est forcé.
  \item[Filtre] préréglages mesurés sur de vraies puces, ou courbe personnalisée.
  \item[Sortie 2SID/3SID] stéréo (puces réparties à gauche, à droite, au centre) ou mono.
  \item[Voie des samples (\$D418)] ajoute à l'oscilloscope une voie par puce montrant les échantillons joués par le
        registre de volume ; les trois voix ne les contiennent alors plus.
\end{description}

\section{SNES}
\capture{settings-spc}{Réglages SNES}
\begin{description}
  \item[Moteur d'émulation] ares ou snes\_spc, appliqué au morceau suivant.
\end{description}

\section{FM}
\capturedouble{settings-ym2612-psg}{settings-ym3812}{Réglages OPN2}{réglages OPL}
Pour chaque famille : cœur d'émulation, préréglage de puce, échelle (« ladder ») du DAC, filtres passe-bas et
passe-haut de sortie, largeur stéréo.

\section{AY et Z80}
\capture{settings-ay}{Réglages AY}
\begin{description}
  \item[Fréquence d'horloge] doit correspondre à la machine d'origine : une horloge fausse décale toutes les notes.
  \item[Type de puce] AY-3-8910 ou YM2149 : le YM a une enveloppe plus fine et des volumes légèrement différents.
  \item[Disposition stéréo] répartition des trois voix (ABC, ACB, mono\dots), selon les habitudes de chaque
        plateforme.
  \item[Interpolation] adoucit le rendu ; désactivée, le son est plus brut, proche de la puce.
\end{description}

\section{Effet des réglages FM}
\begin{description}
  \item[Ladder] les YM2612 d'origine ont un défaut de leur convertisseur qui renforce les sons faibles ; certaines
        bandes-son ont été composées en en tenant compte. Le YM3438 et les Mega Drive tardives ne l'ont pas.
  \item[Filtres de sortie] reproduisent le filtrage analogique de la console ; sans eux, le son est plus brillant
        que sur le matériel.
  \item[Stéréo] largeur de la stéréo : à 0, mono.
\end{description}

\section{Durées et volume}
\begin{description}
  \item[Durée par défaut] utilisée quand ni le format ni une base de données ne donnent la durée.
  \item[Détection du silence] passe au morceau suivant après un silence de la durée choisie.
  \item[Limite de boucles] nombre de répétitions des morceaux qui bouclent.
  \item[Fondus] fondu d'entrée et de sortie.
  \item[Gain] volume général, de 50 à 200~\%.
\end{description}
La liste complète, avec les valeurs par défaut, figure au chapitre~\ref{chap:all-settings}.
